\section{自博弈三阶段框架：从 AlphaGo 轨迹到 LLM 推理扩展}

Brown 把 AlphaGo 和 LLM 的发展轨迹并排比较，形成了一个三阶段模板：
\begin{enumerate}
    \item 在高质量人类数据上预训练（pre-training on human data）；
    \item 放大推理时计算（large-scale inference compute）；
    \item 递归式自改进（recursive self-improvement）。
\end{enumerate}

他用一句非常直接的话概括当前差距：\textit{``In LLMs we don't really have that piece''}，这里 ``that piece'' 指的就是像 AlphaGo 一样的可扩展自博弈自提升。

\begin{importantbox}{三阶段框架的洞察}
真正稀缺的不是 ``模型规模''，而是第三阶段里 \textbf{可证明、可持续、可工程化} 的自改进机制。
\end{importantbox}

这套框架在课堂中的价值，是给研究问题重新排序。传统问法是 ``多 Agent 还能再加什么模块？''；Brown 的问法是 ``为什么在部分环境里 self-play 可以闭环，在另一些环境里会坍缩到脆弱均衡？''。后者更接近根因。

\begin{warningbox}{常见误判}
把 ``LLM 上还没复现 AlphaGo 式提升'' 归因于算力不足，是不充分解释。更常见的根因是任务结构不满足二人零和完美信息假设。
\end{warningbox}

从课程上下文看，这个框架也解释了为什么推理模型（reasoning models）近期进展快，但多 Agent 协作仍显脆弱：前者主要在单体推理链路上吃到了 test-time compute 红利，后者则要额外解决策略相容性与群体分布错配问题。

\subsection{本章小结}
AlphaGo 与 LLM 的类比是有用的，但只能用于定位问题，不能直接迁移结论。决定能否自博弈扩展的，是任务博弈结构而非口号。

